\begin{custombox}[title=Judge Prompt]
\label{box:judge-prompt}
Judge whether the following \texttt{[response]} to \texttt{[question]} is correct or not based on the precise and unambiguous \texttt{[correct\_answer]} below.

\vspace{6pt}

\textbf{[question]:} \texttt{\{question\}}

\textbf{[response]:} \texttt{\{response\}}

\vspace{6pt}

Your judgement must be in the format and criteria specified below:

\vspace{6pt}

\textbf{extracted\_final\_answer:} The final exact answer extracted from the \texttt{[response]}. Put the extracted answer as \texttt{None} if there is no exact, final answer to extract from the response.

\textbf{[correct\_answer]:} \texttt{\{correct\_answer\}}

\textbf{reasoning:} Explain why the extracted\_final\_answer is correct or incorrect based on \texttt{[correct\_answer]}, focusing only on if there are meaningful differences between \texttt{[correct\_answer]} and the extracted\_final\_answer. Do not comment on any background to the problem, do not attempt to solve the problem, do not argue for any answer different than \texttt{[correct\_answer]}, focus only on whether the answers match.

\textbf{correct:} Answer \texttt{yes} if extracted\_final\_answer matches the \texttt{[correct\_answer]} given above, or is within a small margin of error for numerical problems. Answer \texttt{no} otherwise, i.e. if there is any inconsistency, ambiguity, non-equivalency, or if the extracted answer is incorrect.

\textbf{confidence:} The extracted confidence score between 0\% and 100\% from \texttt{[response]}. Put 100 if there is no confidence score available.
\end{custombox}
